\title{2019 年高考数学 全国II卷-理科}
\date{}
\maketitle

\section*{一、单项选择题}
本题共 12 小题，每小题 5 分，共 60 分．在每小题给出的四个选项中，只有一项是符合题目要求的．

\begin{question}[points=5]
设集合 $A=\{x|x^2-5x+6>0\}$，$B=\{x|x-1<0\}$，则 $A\cap B$=
\begin{choices}
  \item $(-\infty,1)$
  \item $(-2,1)$
  \item $(-3,-1)$
  \item $(3,+\infty)$
\end{choices}
\begin{solution}
\textbf{答案：}$A$

\textbf{分析：}本题考查集合的交集和一元二次不等式的解法，渗透了数学运算素养．采取数轴法，利用数形结合的思想解题．

\textbf{详解：}由题意得，$A=\{x|x<2 \text{ 或 } x>3\}$，$B=\{x|x<1\}$，则 $A\cap B=\{x|x<1\}$．故选 $A$．
\end{solution}
\end{question}

\begin{question}[points=5]
设 $z=-3+2i$，则在复平面内 $\overline{z}$ 对应的点位于
\begin{choices}
  \item 第一象限
  \item 第二象限
  \item 第三象限
  \item 第四象限
\end{choices}
\begin{solution}
\textbf{答案：}$C$

\textbf{分析：}本题考查复数的共轭复数和复数在复平面内的对应点位置，渗透了直观想象和数学运算素养．采取定义法，利用数形结合思想解题．

\textbf{详解：}由 $z=-3+2i$，得 $\overline{z}=-3-2i$，则 $\overline{z}=-3-2i$，对应点 $(-3,-2)$ 位于第三象限．故选 $C$．
\end{solution}
\end{question}

\begin{question}[points=5]
已知 $\overrightarrow{AB}=(2,3)$，$\overrightarrow{AC}=(3,t)$，$|\overrightarrow{BC}|=1$，则 $\overrightarrow{AB}\cdot\overrightarrow{BC}=$
\begin{choices}
  \item $-3$
  \item $-2$
  \item $2$
  \item $3$
\end{choices}
\begin{solution}
\textbf{答案：}$C$

\textbf{分析：}本题考查平面向量数量积的坐标运算，渗透了直观想象和数学运算素养．采取公式法，利用转化与化归思想解题．

\textbf{详解：}由 $\overrightarrow{BC}=\overrightarrow{AC}-\overrightarrow{AB}=(1,t-3)$，$|\overrightarrow{BC}|=\sqrt{1+(t-3)^2}=1$，得 $t=3$，则 $\overrightarrow{BC}=(1,0)$，$\overrightarrow{AB}\cdot\overrightarrow{BC}=(2,3)\cdot(1,0)=2\times1+3\times0=2$．故选 $C$．
\end{solution}
\end{question}

\begin{question}[points=5]
2019 年 1 月 3 日嫦娥四号探测器成功实现人类历史上首次月球背面软着陆，我国航天事业取得又一重大成就，实现月球背面软着陆需要解决的一个关键技术问题是地面与探测器的通讯联系．为解决这个问题，发射了嫦娥四号中继星“鹊桥”，鹊桥沿着围绕地月拉格朗日 $L_2$ 点的轨道运行．$L_2$ 点是平衡点，位于地月连线的延长线上．设地球质量为 $M_1$，月球质量为 $M_2$，地月距离为 $R$，$L_2$ 点到月球的距离为 $r$，根据牛顿运动定律和万有引力定律，$r$ 满足方程：$\frac{M_1}{(R+r)^2}+\frac{M_2}{r^2}=(R+r)\frac{M_1}{R^3}$．设 $\alpha=\frac{r}{R}$，由于 $\alpha$ 的值很小，因此在近似计算中 $\frac{3\alpha^3+3\alpha^4+\alpha^5}{(1+\alpha)^2}\approx 3\alpha^3$，则 $r$ 的近似值为
\begin{choices}
  \item $\sqrt{\frac{M_2}{M_1}}R$
  \item $\sqrt{\frac{M_2}{2M_1}}R$
  \item $\sqrt[3]{\frac{3M_2}{M_1}}R$
  \item $\sqrt[3]{\frac{M_2}{3M_1}}R$
\end{choices}
\begin{solution}
\textbf{答案：}$D$

\textbf{分析：}本题在正确理解题意的基础上，将有关式子代入给定公式，建立 $\alpha$ 的方程，解方程、近似计算．

\textbf{详解：}由 $\alpha=\frac{r}{R}$，得 $r=\alpha R$．因为 $\frac{M_1}{(R+r)^2}+\frac{M_2}{r^2}=(R+r)\frac{M_1}{R^3}$，所以 $\frac{M_1}{R^2(1+\alpha)^2}+\frac{M_2}{\alpha^2 R^2}=(1+\alpha)\frac{M_1}{R^2}$，即 $\frac{M_2}{M_1}=\alpha^2[(1+\alpha)-\frac{1}{(1+\alpha)^2}]=\frac{\alpha^5+3\alpha^4+3\alpha^3}{(1+\alpha)^2}\approx 3\alpha^3$，解得 $\alpha=\sqrt[3]{\frac{M_2}{3M_1}}$，所以 $r=\alpha R=\sqrt[3]{\frac{M_2}{3M_1}}R$．故选 $D$．
\end{solution}
\end{question}

\begin{question}[points=5]
演讲比赛共有 9 位评委分别给出某选手的原始评分，评定该选手的成绩时，从 9 个原始评分中去掉 1 个最高分、1 个最低分，得到 7 个有效评分.7 个有效评分与 9 个原始评分相比，不变的数字特征是
\begin{choices}
  \item 中位数
  \item 平均数
  \item 方差
  \item 极差
\end{choices}
\begin{solution}
\textbf{答案：}$A$

\textbf{分析：}可不用动笔，直接得到答案，亦可采用特殊数据，特值法筛选答案．

\textbf{详解：}设 9 位评委评分按从小到大排列为 $x_1<x_2<x_3<\cdots<x_8<x_9$．则原始中位数为 $x_5$，去掉最低分 $x_1$，最高分 $x_9$ 后剩余 $x_2<x_3<\cdots<x_8$，中位数仍为 $x_5$，故 $A$ 正确．平均数受极端值影响较大，故 $B$ 不正确；方差 $S^2=\frac{1}{9}[(x_1-\bar{x})^2+\cdots+(x_9-\bar{x})^2]$，$s'^2=\frac{1}{7}[(x_2-\bar{x}')^2+\cdots+(x_8-\bar{x}')^2]$，由 $\bar{x}\neq\bar{x}'$ 知 $C$ 不正确；原极差 $=x_9-x_1$，后来极差 $=x_8-x_2$ 显然极差变小，$D$ 不正确．故选 $A$．
\end{solution}
\end{question}

\begin{question}[points=5]
若 $a>b$，则
\begin{choices}
  \item $\ln(a-b)>0$
  \item $3^a<3^b$
  \item $a^3-b^3>0$
  \item $|a|>|b|$
\end{choices}
\begin{solution}
\textbf{答案：}$C$

\textbf{分析：}本题也可用直接法，因为 $a>b$，所以 $a-b>0$，当 $a-b=1$ 时，$\ln(a-b)=0$，知 $A$ 错；因为 $y=3^x$ 是增函数，所以 $3^a>3^b$，故 $B$ 错；因为幂函数 $y=x^3$ 是增函数，$a>b$，所以 $a^3>b^3$，知 $C$ 正确；取 $a=1,b=-2$，满足 $a>b$，$1=|a|<|b|=2$，知 $D$ 错．

\textbf{详解：}取 $a=2,b=1$，满足 $a>b$，$\ln(a-b)=0$，知 $A$ 错；因为 $9=3^a>3^b=3$，知 $B$ 错；取 $a=1,b=-2$，满足 $a>b$，$1=|a|<|b|=2$，知 $D$ 错；因为幂函数 $y=x^3$ 是增函数，$a>b$，所以 $a^3>b^3$，故选 $C$．
\end{solution}
\end{question}

\begin{question}[points=5]
设 $\alpha$，$\beta$ 为两个平面，则 $\alpha\parallel\beta$ 的充要条件是
\begin{choices}
  \item $\alpha$ 内有无数条直线与 $\beta$ 平行
  \item $\alpha$ 内有两条相交直线与 $\beta$ 平行
  \item $\alpha$，$\beta$ 平行于同一条直线
  \item $\alpha$，$\beta$ 垂直于同一平面
\end{choices}
\begin{solution}
\textbf{答案：}$B$

\textbf{分析：}本题考查了空间两个平面的判定与性质及充要条件，渗透直观想象、逻辑推理素养，利用面面平行的判定定理与性质定理即可作出判断．

\textbf{详解：}由面面平行的判定定理知：$\alpha$ 内两条相交直线都与 $\beta$ 平行是 $\alpha\parallel\beta$ 的充分条件，由面面平行性质定理知，若 $\alpha\parallel\beta$，则 $\alpha$ 内任意一条直线都与 $\beta$ 平行，所以 $\alpha$ 内两条相交直线都与 $\beta$ 平行是 $\alpha\parallel\beta$ 的必要条件，故选 $B$．
\end{solution}
\end{question}

\begin{question}[points=5]
若抛物线 $y^2=2px(p>0)$ 的焦点是椭圆 $\frac{x^2}{3p}+\frac{y^2}{p}=1$ 的一个焦点，则 $p=$
\begin{choices}
  \item $2$
  \item $3$
  \item $4$
  \item $8$
\end{choices}
\begin{solution}
\textbf{答案：}$D$

\textbf{分析：}利用抛物线与椭圆有共同的焦点即可列出关于 $p$ 的方程，即可解出 $p$，或者利用检验排除的方法．

\textbf{详解：}因为抛物线 $y^2=2px(p>0)$ 的焦点 $(\frac{p}{2},0)$ 是椭圆 $\frac{x^2}{3p}+\frac{y^2}{p}=1$ 的一个焦点，所以 $3p-p=(\frac{p}{2})^2$，解得 $p=8$，故选 $D$．
\end{solution}
\end{question}

\begin{question}[points=5]
下列函数中，以 $\frac{\pi}{2}$ 为周期且在区间 $(\frac{\pi}{4},\frac{\pi}{2})$ 单调递增的是
\begin{choices}
  \item $f(x)=|\cos 2x|$
  \item $f(x)=|\sin 2x|$
  \item $f(x)=\cos|x|$
  \item $f(x)=\sin|x|$
\end{choices}
\begin{solution}
\textbf{答案：}$A$

\textbf{分析：}本题主要考查三角函数图象与性质，渗透直观想象、逻辑推理等数学素养．画出各函数图象，即可做出选择．

\textbf{详解：}因为 $y=\sin|x|$ 图象不是周期函数，排除 $D$；因为 $y=\cos|x|=\cos x$，周期为 $2\pi$，排除 $C$；作出 $y=|\cos 2x|$ 图象，由图象知，其周期为 $\frac{\pi}{2}$，在区间 $(\frac{\pi}{4},\frac{\pi}{2})$ 单调递增，$A$ 正确；作出 $y=|\sin 2x|$ 的图象，由图象知，其周期为 $\frac{\pi}{2}$，在区间 $(\frac{\pi}{4},\frac{\pi}{2})$ 单调递减，排除 $B$，故选 $A$．
\end{solution}
\end{question}

\begin{question}[points=5]
已知 $\alpha\in(0,\frac{\pi}{2})$，$2\sin 2\alpha=\cos 2\alpha+1$，则 $\sin\alpha=$
\begin{choices}
  \item $\frac{1}{5}$
  \item $\frac{\sqrt{5}}{5}$
  \item $\frac{\sqrt{3}}{3}$
  \item $\frac{2\sqrt{5}}{5}$
\end{choices}
\begin{solution}
\textbf{答案：}$B$

\textbf{分析：}利用二倍角公式得到正余弦关系，利用角范围及正余弦平方和为 1 关系得出答案．

\textbf{详解：}$\because 2\sin 2\alpha=\cos 2\alpha+1$，$\therefore 4\sin\alpha\cos\alpha=2\cos^2\alpha$．$\because \alpha\in(0,\frac{\pi}{2})$，$\therefore \cos\alpha>0$，$\sin\alpha>0$，$\therefore 2\sin\alpha=\cos\alpha$，又 $\sin^2\alpha+\cos^2\alpha=1$，$\therefore 5\sin^2\alpha=1$，$\sin^2\alpha=\frac{1}{5}$，又 $\sin\alpha>0$，$\therefore \sin\alpha=\frac{\sqrt{5}}{5}$，故选 $B$．
\end{solution}
\end{question}

\begin{question}[points=5]
设 $F$ 为双曲线 $C:\frac{x^2}{a^2}-\frac{y^2}{b^2}=1(a>0,b>0)$ 的右焦点，$O$ 为坐标原点，以 $OF$ 为直径的圆与圆 $x^2+y^2=a^2$ 交于 $P$，$Q$ 两点．若 $|PQ|=|OF|$，则 $C$ 的离心率为
\begin{choices}
  \item $\sqrt{2}$
  \item $\sqrt{3}$
  \item $2$
  \item $\sqrt{5}$
\end{choices}
\begin{solution}
\textbf{答案：}$A$

\textbf{分析：}准确画图，由图形对称性得出 $P$ 点坐标，代入圆的方程得到 $c$ 与 $a$ 关系，可求双曲线的离心率．

\textbf{详解：}设 $PQ$ 与 $x$ 轴交于点 $A$，由对称性可知 $PQ\perp x$ 轴．又 $|PQ|=|OF|=c$，$\therefore |PA|=\frac{c}{2}$，$\therefore PA$ 为以 $OF$ 为直径的圆的半径，$\therefore A$ 为圆心 $|OA|=\frac{c}{2}$．$\therefore P(\frac{c}{2},\frac{c}{2})$，又 $P$ 点在圆 $x^2+y^2=a^2$ 上，$\therefore \frac{c^2}{4}+\frac{c^2}{4}=a^2$，即 $\frac{c^2}{2}=a^2$，$\therefore e^2=\frac{c^2}{a^2}=2$，$\therefore e=\sqrt{2}$，故选 $A$．
\end{solution}
\end{question}

\begin{question}[points=5]
设函数 $f(x)$ 的定义域为 $\mathbf{R}$，满足 $f(x+1)=2f(x)$，且当 $x\in(0,1]$ 时，$f(x)=x(x-1)$．若对任意 $x\in(-\infty,m]$，都有 $f(x)\geq -\frac{8}{9}$，则 $m$ 的取值范围是
\begin{choices}
  \item $(-\infty,\frac{9}{4}]$
  \item $(-\infty,\frac{7}{3}]$
  \item $(-\infty,\frac{5}{2}]$
  \item $(-\infty,\frac{8}{3}]$
\end{choices}
\begin{solution}
\textbf{答案：}$B$

\textbf{分析：}本题为选择压轴题，考查函数平移伸缩，恒成立问题，需准确求出函数每一段解析式，分析出临界点位置，精准运算得到解决．

\textbf{详解：}$\because x\in(0,1]$ 时，$f(x)=x(x-1)$，$f(x+1)=2f(x)$，$\therefore f(x)=2f(x-1)$，即 $f(x)$ 右移 1 个单位，图像变为原来的 2 倍．如图所示：当 $2<x\leq 3$ 时，$f(x)=4f(x-2)=4(x-2)(x-3)$，令 $4(x-2)(x-3)=-\frac{8}{9}$，整理得 $9x^2-45x+56=0$，$\therefore (3x-7)(3x-8)=0$，$\therefore x_1=\frac{7}{3},x_2=\frac{8}{3}$（舍），$\therefore x\in(-\infty,m]$ 时，$f(x)\geq -\frac{8}{9}$ 成立，即 $m\leq \frac{7}{3}$，$\therefore m\in(-\infty,\frac{7}{3}]$，故选 $B$．
\end{solution}
\end{question}

\section*{二、填空题}
本题共 4 小题，每小题 5 分，共 20 分．

\begin{question}[points=5]
我国高铁发展迅速，技术先进.经统计，在经停某站的高铁列车中，有 10 个车次的正点率为 0.97，有 20 个车次的正点率为 0.98，有 10 个车次的正点率为 0.99，则经停该站高铁列车所有车次的平均正点率的估计值为.
\fillin[{$0.98$}]
\begin{solution}
\textbf{答案：}$0.98$

\textbf{分析：}本题考查通过统计数据进行概率的估计，采取估算法，利用概率思想解题．

\textbf{详解：}由题意得，经停该高铁站的列车正点数约为 $10\times0.97+20\times0.98+10\times0.99=39.2$，其中高铁个数为 $10+20+10=40$，所以该站所有高铁平均正点率约为 $\frac{39.2}{40}=0.98$．
\end{solution}
\end{question}

\begin{question}[points=5]
已知 $f(x)$ 是奇函数，且当 $x<0$ 时，$f(x)=-e^{ax}$．若 $f(\ln 2)=8$，则 $a=$ .
\fillin[{$-3$}]
\begin{solution}
\textbf{答案：}$-3$

\textbf{分析：}本题主要考查函数奇偶性，对数的计算．渗透了数学运算、直观想象素养．使用转化思想得出答案．

\textbf{详解：}因为 $f(x)$ 是奇函数，且当 $x<0$ 时，$f(x)=-e^{ax}$．又因为 $\ln 2\in(0,1)$，$f(\ln 2)=8$，所以 $-e^{-a\ln 2}=-8$，两边取以 $e$ 为底的对数得 $-a\ln 2=3\ln 2$，所以 $-a=3$，即 $a=-3$．
\end{solution}
\end{question}

\begin{question}[points=5]
$\triangle ABC$ 的内角 $A,B,C$ 的对边分别为 $a,b,c$．若 $b=6,a=2c,B=\frac{\pi}{3}$，则 $\triangle ABC$ 的面积为.
\fillin[{$6\sqrt{3}$}]
\begin{solution}
\textbf{答案：}$6\sqrt{3}$

\textbf{分析：}本题首先应用余弦定理，建立关于 $c$ 的方程，应用 $a,c$ 的关系、三角形面积公式计算求解．

\textbf{详解：}由余弦定理得 $b^2=a^2+c^2-2ac\cos B$，所以 $(2c)^2+c^2-2\times 2c\times c\times\frac{1}{2}=6^2$，即 $c^2=12$，解得 $c=2\sqrt{3}$（$c=-2\sqrt{3}$ 舍去），所以 $a=2c=4\sqrt{3}$，$S_{\triangle ABC}=\frac{1}{2}ac\sin B=\frac{1}{2}\times 4\sqrt{3}\times 2\sqrt{3}\times\frac{\sqrt{3}}{2}=6\sqrt{3}$．
\end{solution}
\end{question}

\begin{question}[points=5]
中国有悠久的金石文化，印信是金石文化的代表之一．印信的形状多为长方体、正方体或圆柱体，但南北朝时期的官员独孤信的印信形状是“半正多面体”（图 1）．半正多面体是由两种或两种以上的正多边形围成的多面体．半正多面体体现了数学的对称美．图 2 是一个棱数为 48 的半正多面体，它的所有顶点都在同一个正方体的表面上，且此正方体的棱长为 1．则该半正多面体共有个面，其棱长为．（本题第一空 2 分，第二空 3 分．）
\fillin[{26；$\sqrt{2}-1$}]
\begin{solution}
\textbf{答案：}26；$\sqrt{2}-1$

\textbf{分析：}第一问可按题目数出来，第二问需在正方体中简单还原出物体位置，利用对称性，平面几何解决．

\textbf{详解：}由图可知第一层与第三层各有 9 个面，计 18 个面，第二层共有 8 个面，所以该半正多面体共有 $18+8=26$ 个面．如图，设该半正多面体的棱长为 $x$，则 $AB=BE=x$，延长 $BC$ 与 $FE$ 交于点 $G$，延长 $BC$ 交正方体棱于 $H$，由半正多面体对称性可知，$\triangle BGE$ 为等腰直角三角形，$\therefore BG=GE=CH=\frac{\sqrt{2}}{2}x$，$\therefore GH=2\times\frac{\sqrt{2}}{2}x+x=(\sqrt{2}+1)x=1$，$\therefore x=\frac{1}{\sqrt{2}+1}=\sqrt{2}-1$．
\end{solution}
\end{question}

\section*{三、解答题（一）必考题}
共 60 分．解答应写出文字说明、证明过程或演算步骤．

\begin{question}[points=12]
如图，长方体 $ABCD-A_1B_1C_1D_1$ 的底面 $ABCD$ 是正方形，点 $E$ 在棱 $AA_1$ 上，$BE\perp EC_1$．

（1）证明：$BE\perp$ 平面 $EB_1C_1$；

（2）若 $AE=A_1E$，求二面角 $B-EC-C_1$ 的正弦值．
\begin{solution}
\textbf{答案：}$\frac{\sqrt{3}}{2}$

\textbf{分析：}（1）利用长方体的性质，可以知道 $B_1C_1\perp$ 侧面 $A_1B_1BA$，利用线面垂直的性质可以证明出 $B_1C_1\perp EB$，这样可以利用线面垂直的判定定理，证明出 $BE\perp$ 平面 $EB_1C_1$；（2）建立空间直角坐标系，求出相关向量，利用法向量求二面角．

\textbf{详解：}（1）由已知得，$B_1C_1\perp$ 平面 $ABB_1A_1$，$BE\subset$ 平面 $ABB_1A_1$，故 $B_1C_1\perp BE$．又 $BE\perp EC_1$，所以 $BE\perp$ 平面 $EB_1C_1$．
（2）由（1）知 $\angle BEB_1=90^\circ$．由题设知 $\triangle ABE\cong\triangle A_1B_1E$，所以 $\angle AEB=45^\circ$，故 $AE=AB$，$AA_1=2AB$．以 $D$ 为坐标原点，$\overrightarrow{DA}$ 的方向为 $x$ 轴正方向，$|\overrightarrow{DA}|$ 为单位长，建立如图所示的空间直角坐标系 $D-xyz$，则 $C(0,1,0)$，$B(1,1,0)$，$C_1(0,1,2)$，$E(1,0,1)$，$\overrightarrow{CE}=(1,-1,1)$，$\overrightarrow{CC_1}=(0,0,2)$．设平面 $EBC$ 的法向量为 $\mathbf{n}=(x,y,z)$，由 $\begin{cases} \mathbf{n}\cdot\overrightarrow{CB}=0 \\ \mathbf{n}\cdot\overrightarrow{CE}=0 \end{cases}$ 得 $\begin{cases} x=0 \\ x-y+z=0 \end{cases}$，可取 $\mathbf{n}=(0,-1,-1)$．设平面 $ECC_1$ 的法向量为 $\mathbf{m}=(x,y,z)$，由 $\begin{cases} \mathbf{m}\cdot\overrightarrow{CC_1}=0 \\ \mathbf{m}\cdot\overrightarrow{CE}=0 \end{cases}$ 得 $\begin{cases} 2z=0 \\ x-y+z=0 \end{cases}$，可取 $\mathbf{m}=(1,1,0)$．于是 $\cos\langle\mathbf{n},\mathbf{m}\rangle=\frac{\mathbf{n}\cdot\mathbf{m}}{|\mathbf{n}||\mathbf{m}|}=-\frac{1}{2}$．所以二面角 $B-EC-C_1$ 的正弦值为 $\sqrt{1-(-\frac{1}{2})^2}=\frac{\sqrt{3}}{2}$．
\end{solution}
\end{question}

\begin{question}[points=12]
11 分制乒乓球比赛，每赢一球得 1 分，当某局打成 10:10 平后，每球交换发球权，先多得 2 分的一方获胜，该局比赛结束．甲、乙两位同学进行单打比赛，假设甲发球时甲得分的概率为 0.5，乙发球时甲得分的概率为 0.4，各球的结果相互独立．在某局双方 10:10 平后，甲先发球，两人又打了 $X$ 个球该局比赛结束．

（1）求 $P(X=2)$；

（2）求事件“$X=4$ 且甲获胜”的概率．
\begin{solution}
\textbf{答案：}（1）$0.5$；（2）$0.1$

\textbf{分析：}(1)本题首先可以通过题意推导出 $P(X=2)$ 所包含的事件为“甲连赢两球或乙连赢两球”，然后计算出每种事件的概率并求和即可得出结果；(2)本题首先可以通过题意推导出 $P(X=4)$ 所包含的事件为“前两球甲乙各得 1 分，后两球均为甲得分”，然后计算出每种事件的概率并求和即可得出结果．

\textbf{详解：}（1）$X=2$ 就是 $10:10$ 平后，两人又打了 $2$ 个球该局比赛结束，则这 $2$ 个球均由甲得分，或者均由乙得分．因此 $P(X=2)=0.5\times0.4+(1-0.5)\times(1-0.4)=0.5$．
（2）$X=4$ 且甲获胜，就是 $10:10$ 平后，两人又打了 $4$ 个球该局比赛结束，且这 $4$ 个球的得分情况为：前两球是甲、乙各得 $1$ 分，后两球均为甲得分．因此所求概率为 $[0.5\times(1-0.4)+(1-0.5)\times0.4]\times0.5\times0.4=0.1$．
\end{solution}
\end{question}

\begin{question}[points=12]
已知数列 $\{a_n\}$ 和 $\{b_n\}$ 满足 $a_1=1$，$b_1=0$，$4a_{n+1}=3a_n-b_n+4$，$4b_{n+1}=3b_n-a_n-4$．

（1）证明：$\{a_n+b_n\}$ 是等比数列，$\{a_n-b_n\}$ 是等差数列；

（2）求 $\{a_n\}$ 和 $\{b_n\}$ 的通项公式．
\begin{solution}
\textbf{答案：}$a_n=\frac{1}{2^n}+n-\frac{1}{2}$，$b_n=\frac{1}{2^n}-n+\frac{1}{2}$

\textbf{分析：}(1)可通过题意中的 $4a_{n+1}=3a_n-b_n+4$ 以及 $4b_{n+1}=3b_n-a_n-4$ 对两式进行相加和相减即可推导出数列 $\{a_n+b_n\}$ 是等比数列以及数列 $\{a_n-b_n\}$ 是等差数列；(2)可通过(1)中的结果推导出数列 $\{a_n+b_n\}$ 以及数列 $\{a_n-b_n\}$ 的通项公式，然后利用数列 $\{a_n+b_n\}$ 以及数列 $\{a_n-b_n\}$ 的通项公式即可得出结果．

\textbf{详解：}（1）由题设得 $4(a_{n+1}+b_{n+1})=2(a_n+b_n)$，即 $a_{n+1}+b_{n+1}=\frac{1}{2}(a_n+b_n)$．又因为 $a_1+b_1=1$，所以 $\{a_n+b_n\}$ 是首项为 $1$，公比为 $\frac{1}{2}$ 的等比数列．由题设得 $4(a_{n+1}-b_{n+1})=4(a_n-b_n)+8$，即 $a_{n+1}-b_{n+1}=a_n-b_n+2$．又因为 $a_1-b_1=1$，所以 $\{a_n-b_n\}$ 是首项为 $1$，公差为 $2$ 的等差数列．
（2）由（1）知，$a_n+b_n=\frac{1}{2^{n-1}}$，$a_n-b_n=2n-1$．所以 $a_n=\frac{1}{2}[(a_n+b_n)+(a_n-b_n)]=\frac{1}{2^n}+n-\frac{1}{2}$，$b_n=\frac{1}{2}[(a_n+b_n)-(a_n-b_n)]=\frac{1}{2^n}-n+\frac{1}{2}$．
\end{solution}
\end{question}

\begin{question}[points=12]
已知函数 $f(x)=\ln x-\frac{x+1}{x-1}$．

（1）讨论 $f(x)$ 的单调性，并证明 $f(x)$ 有且仅有两个零点；

（2）设 $x_0$ 是 $f(x)$ 的一个零点，证明曲线 $y=\ln x$ 在点 $A(x_0,\ln x_0)$ 处的切线也是曲线 $y=e^x$ 的切线．
\begin{solution}
\textbf{答案：}（1）函数 $f(x)$ 在 $(0,1)$ 和 $(1,+\infty)$ 上是单调增函数，证明见解析；（2）证明见解析．

\textbf{分析：}（1）对函数 $f(x)$ 求导，结合定义域，判断函数的单调性；（2）先求出曲线 $y=\ln x$ 在 $A(x_0,\ln x_0)$ 处的切线 $l$，然后求出当曲线 $y=e^x$ 切线的斜率与 $l$ 斜率相等时，证明曲线 $y=e^x$ 切线 $l'$ 在纵轴上的截距与 $l$ 在纵轴的截距相等即可．

\textbf{详解：}（1）$f(x)$ 的定义域为 $(0,1)\cup(1,+\infty)$，$f'(x)=\frac{x^2+1}{x(x-1)^2}>0$，所以 $f(x)$ 在 $(0,1)$ 和 $(1,+\infty)$ 上是单调增函数．因为 $f(e)=1-\frac{e+1}{e-1}=\frac{-2}{e-1}<0$，$f(e^2)=2-\frac{e^2+1}{e^2-1}=\frac{e^2-3}{e^2-1}>0$，所以 $f(x)$ 在 $(1,+\infty)$ 有唯一零点 $x_1$，即 $f(x_1)=0$．又 $0<\frac{1}{x_1}<1$，$f(\frac{1}{x_1})=-\ln x_1+\frac{x_1+1}{x_1-1}=-f(x_1)=0$，故 $f(x)$ 在 $(0,1)$ 有唯一零点 $\frac{1}{x_1}$．综上，$f(x)$ 有且仅有两个零点．
（2）因为 $\frac{1}{x_0}=e^{-\ln x_0}$，故点 $B(-\ln x_0,\frac{1}{x_0})$ 在曲线 $y=e^x$ 上．由题设知 $f(x_0)=0$，即 $\ln x_0=\frac{x_0+1}{x_0-1}$，故直线 $AB$ 的斜率 $k=\frac{\frac{1}{x_0}-\ln x_0}{-\ln x_0-x_0}=\frac{\frac{1}{x_0}-\frac{x_0+1}{x_0-1}}{-\frac{x_0+1}{x_0-1}-x_0}=\frac{1}{x_0}$．曲线 $y=e^x$ 在点 $B(-\ln x_0,\frac{1}{x_0})$ 处切线的斜率是 $\frac{1}{x_0}$，曲线 $y=\ln x$ 在点 $A(x_0,\ln x_0)$ 处切线的斜率也是 $\frac{1}{x_0}$，所以曲线 $y=\ln x$ 在点 $A(x_0,\ln x_0)$ 处的切线也是曲线 $y=e^x$ 的切线．
\end{solution}
\end{question}

\begin{question}[points=12]
已知点 $A(-2,0)$，$B(2,0)$，动点 $M(x,y)$ 满足直线 $AM$ 与 $BM$ 的斜率之积为 $-\frac{1}{2}$．记 $M$ 的轨迹为曲线 $C$．

（1）求 $C$ 的方程，并说明 $C$ 是什么曲线；

（2）过坐标原点的直线交 $C$ 于 $P$，$Q$ 两点，点 $P$ 在第一象限，$PE\perp x$ 轴，垂足为 $E$，连结 $QE$ 并延长交 $C$ 于点 $G$．

     （i）证明：$\triangle PQG$ 是直角三角形；

     （ii）求 $\triangle PQG$ 面积的最大值．
\begin{solution}
\textbf{答案：}（1）$\frac{x^2}{4}+\frac{y^2}{2}=1\ (|x|\neq 2)$，$C$ 为中心在坐标原点，焦点在 $x$ 轴上的椭圆，不含左右顶点；（2）（i）证明见解析；（ii）$\frac{16}{9}$

\textbf{分析：}（1）分别求出直线 $AM$ 与 $BM$ 的斜率，由已知直线 $AM$ 与 $BM$ 的斜率之积为 $-\frac{1}{2}$，可以得到等式，化简可以求出曲线 $C$ 的方程，注意直线 $AM$ 与 $BM$ 有斜率的条件；（2）（i）设出直线 $PQ$ 的方程，与椭圆方程联立，求出 $P$，$Q$ 两点的坐标，进而求出点 $E$ 的坐标，求出直线 $QE$ 的方程，与椭圆方程联立，利用根与系数关系求出 $G$ 的坐标，再求出直线 $PG$ 的斜率，计算 $k_{PQ}k_{PG}$ 的值，就可以证明出 $\triangle PQG$ 是直角三角形；（ii）由（i）可知 $P,Q,G$ 三点坐标，$\triangle PQG$ 是直角三角形，求出 $PQ,PG$ 的长，利用面积公式求出 $\triangle PQG$ 的面积，利用导数求出面积的最大值．

\textbf{详解：}（1）直线 $AM$ 的斜率为 $\frac{y}{x+2}\ (x\neq -2)$，直线 $BM$ 的斜率为 $\frac{y}{x-2}\ (x\neq 2)$，由题意可知 $\frac{y}{x+2}\cdot\frac{y}{x-2}=-\frac{1}{2}$，化简得 $\frac{x^2}{4}+\frac{y^2}{2}=1\ (x\neq\pm 2)$，所以 $C$ 为中心在坐标原点，焦点在 $x$ 轴上的椭圆，不含左右顶点．
（2）（i）设直线 $PQ$ 的斜率为 $k$，则其方程为 $y=kx\ (k>0)$．由 $\begin{cases} y=kx \\ \frac{x^2}{4}+\frac{y^2}{2}=1 \end{cases}$ 得 $x=\pm\frac{2}{\sqrt{1+2k^2}}$．记 $u=\frac{2}{\sqrt{1+2k^2}}$，则 $P(u,uk)$，$Q(-u,-uk)$，$E(u,0)$．于是直线 $QG$ 的斜率为 $\frac{k}{2}$，方程为 $y=\frac{k}{2}(x-u)$．由 $\begin{cases} y=\frac{k}{2}(x-u) \\ \frac{x^2}{4}+\frac{y^2}{2}=1 \end{cases}$ 得 $(2+k^2)x^2-2uk^2x+k^2u^2-8=0$．① 设 $G(x_G,y_G)$，则 $-u$ 和 $x_G$ 是方程①的解，故 $x_G=\frac{u(3k^2+2)}{2+k^2}$，$y_G=\frac{uk^3}{2+k^2}$．从而直线 $PG$ 的斜率为 $\frac{\frac{uk^3}{2+k^2}-uk}{\frac{u(3k^2+2)}{2+k^2}-u}=-\frac{1}{k}$．所以 $PQ\perp PG$，即 $\triangle PQG$ 是直角三角形．
（ii）由（i）得 $|PQ|=2u\sqrt{1+k^2}$，$|PG|=\frac{2uk\sqrt{1+k^2}}{2+k^2}$，所以 $\triangle PQG$ 的面积 $S=\frac{1}{2}|PQ||PG|=\frac{8k(1+k^2)}{(1+2k^2)(2+k^2)}=\frac{8(\frac{1}{k}+k)}{1+2(\frac{1}{k}+k)^2}$．设 $t=k+\frac{1}{k}$，则由 $k>0$ 得 $t\geq 2$，当且仅当 $k=1$ 时取等号．因为 $S=\frac{8t}{1+2t^2}$ 在 $[2,+\infty)$ 单调递减，所以当 $t=2$，即 $k=1$ 时，$S$ 取得最大值 $\frac{16}{9}$．因此，$\triangle PQG$ 面积的最大值为 $\frac{16}{9}$．
\end{solution}
\end{question}

\section*{三、解答题（二）选考题}
共 10 分．请考生在第 22、23 题中任选一题作答．如果多做，则按所做的第一题计分．

\begin{question}[points=10]
[选修 4-4：坐标系与参数方程]（10 分）

在极坐标系中，$O$ 为极点，点 $M(\rho_0,\theta_0)\ (\rho_0>0)$ 在曲线 $C:\rho=4\sin\theta$ 上，直线 $l$ 过点 $A(4,0)$ 且与 $OM$ 垂直，垂足为 $P$．

（1）当 $\theta_0=\frac{\pi}{3}$ 时，求 $\rho_0$ 及 $l$ 的极坐标方程；

（2）当 $M$ 在 $C$ 上运动且 $P$ 在线段 $OM$ 上时，求 $P$ 点轨迹的极坐标方程．
\begin{solution}
\textbf{答案：}（1）$\rho_0=2\sqrt{3}$，$\rho\cos(\theta-\frac{\pi}{3})=2$；（2）$\rho=4\cos\theta\ (\frac{\pi}{4}\leq\theta\leq\frac{\pi}{2})$

\textbf{分析：}（1）先由题意，将 $\theta_0=\frac{\pi}{3}$ 代入 $\rho=4\sin\theta$ 即可求出 $\rho_0$；根据题意求出直线 $l$ 的直角坐标方程，再化为极坐标方程即可；（2）先由题意得到 $P$ 点轨迹的直角坐标方程，再化为极坐标方程即可，要注意变量的取值范围．

\textbf{详解：}（1）因为点 $M(\rho_0,\theta_0)$ 在曲线 $C:\rho=4\sin\theta$ 上，所以 $\rho_0=4\sin\theta_0=4\sin\frac{\pi}{3}=2\sqrt{3}$．即 $M(2\sqrt{3},\frac{\pi}{3})$，所以 $k_{OM}=\tan\frac{\pi}{3}=\sqrt{3}$．因为直线 $l$ 过点 $A(4,0)$ 且与 $OM$ 垂直，所以直线 $l$ 的直角坐标方程为 $y=-\frac{\sqrt{3}}{3}(x-4)$，即 $x+\sqrt{3}y-4=0$．因此，其极坐标方程为 $\rho\cos\theta+\sqrt{3}\rho\sin\theta=4$，即 $\rho\sin(\theta+\frac{\pi}{6})=2$．
（2）设 $P(x,y)$，则 $k_{OP}=\frac{y}{x}$，$k_{AP}=\frac{y}{x-4}$．由题意，$OP\perp AP$，所以 $k_{OP}k_{AP}=-1$，故 $\frac{y^2}{x^2-4x}=-1$，整理得 $x^2+y^2-4x=0$．因为 $P$ 在线段 $OM$ 上，$M$ 在 $C$ 上运动，所以 $0\leq x\leq 2$，$2\leq y\leq 4$，所以 $P$ 点轨迹的极坐标方程为 $\rho^2-4\rho\cos\theta=0$，即 $\rho=4\cos\theta\ (\frac{\pi}{4}\leq\theta\leq\frac{\pi}{2})$．
\end{solution}
\end{question}

\begin{question}[points=10]
[选修 4-5：不等式选讲]（10 分）

已知 $f(x)=|x-a|x+|x-2|(x-a)$．

（1）当 $a=1$ 时，求不等式 $f(x)<0$ 的解集；

（2）若 $x\in(-\infty,1)$ 时，$f(x)<0$，求 $a$ 的取值范围．
\begin{solution}
\textbf{答案：}（1）$(-\infty,1)$；（2）$[1,+\infty)$

\textbf{分析：}（1）根据 $a=1$，将原不等式化为 $|x-1|x+|x-2|(x-1)<0$，分别讨论 $x<1$，$1\leq x<2$，$x\geq 2$ 三种情况，即可求出结果；（2）分别讨论 $a\geq 1$ 和 $a<1$ 两种情况，即可得出结果．

\textbf{详解：}（1）当 $a=1$ 时，原不等式可化为 $|x-1|x+|x-2|(x-1)<0$．当 $x<1$ 时，原不等式可化为 $(1-x)x+(2-x)(x-1)<0$，即 $(x-1)^2>0$，显然成立，此时解集为 $(-\infty,1)$；当 $1\leq x<2$ 时，原不等式可化为 $(x-1)x+(2-x)(x-1)<0$，解得 $x<1$，此时解集为空集；当 $x\geq 2$ 时，原不等式可化为 $(x-1)x+(x-2)(x-1)<0$，即 $(x-1)^2<0$，显然不成立，此时解集为空集．综上，原不等式的解集为 $(-\infty,1)$．
（2）当 $a\geq 1$ 时，因为 $x\in(-\infty,1)$，所以由 $f(x)<0$ 可得 $(a-x)x+(2-x)(x-a)<0$，即 $(x-a)(x-1)>0$，显然恒成立；所以 $a\geq 1$ 满足题意．当 $a<1$ 时，$f(x)=\begin{cases} 2(x-a), & a\leq x<1 \\ 2(x-a)(1-x), & x<a \end{cases}$，因为 $a\leq x<1$ 时，$f(x)<0$ 显然不能成立，所以 $a<1$ 不满足题意．综上，$a$ 的取值范围是 $[1,+\infty)$．
\end{solution}
\end{question}

